\documentclass[9pt,a4paper]{article}

\usepackage{iftex}
\ifPDFTeX
    \usepackage[utf8]{inputenc}
    \usepackage[T5]{fontenc}
    \usepackage[vietnamese]{babel}
    \usepackage{lmodern}
\else
    \usepackage{fontspec}
    \setmainfont{Times New Roman}
\fi
\usepackage{geometry}
\usepackage{longtable}
\usepackage{booktabs}
\usepackage{array}
\usepackage{enumitem}
\usepackage{hyperref}
\usepackage{xcolor}
\usepackage{amsmath}
\usepackage{amssymb}
\usepackage{microtype}
\usepackage[compact]{titlesec}
\titlespacing*{\section}{0pt}{5pt}{2pt}
\titlespacing*{\subsection}{0pt}{3pt}{1pt}
\titleformat*{\section}{\normalsize\bfseries}
\titleformat*{\subsection}{\small\bfseries}
\setlist{noitemsep,topsep=1pt,leftmargin=*}

\geometry{margin=1.1cm}
\hypersetup{colorlinks=true, linkcolor=blue, urlcolor=blue, citecolor=blue}

% Cho phép LaTeX giãn dòng linh hoạt hơn, tránh khoảng cách từ bị kéo quá xa
\emergencystretch=1.5em
\tolerance=800
\hbadness=10000
\setlength{\emergencystretch}{1.5em}

\renewcommand{\arraystretch}{1.0}

% Equation spacing
\setlength{\abovedisplayskip}{4pt}
\setlength{\belowdisplayskip}{4pt}
\setlength{\abovedisplayshortskip}{2pt}
\setlength{\belowdisplayshortskip}{2pt}

\begin{document}

\begin{center}
    \large\textbf{KIẾN TRÚC ĐỊNH THỜI LOGIC PHÂN TÁN (STRICT VS HEURISTIC)} \\
    \small\textit{Tài liệu Tóm tắt Thiết kế --- Bài toán \#112 Log Delay Compensator (2-Page Executive Brief)}
\end{center}

\section{Kiến trúc Tổng thể \& Bài toán}
Hệ thống xử lý luồng log Web Server đến \emph{out-of-order} 24/7, dùng \textbf{Watermark} để quyết định chờ dữ liệu trễ bao lâu trước khi chốt cửa sổ (Tumbling Window 5s trên Event-Time). Deliverable cốt lõi là đường cong đánh đổi \textbf{Data Completeness \% vs Wait Time}. Hệ thống được phân rã thành các phân lớp độc lập, dùng chung một đường dữ liệu (data plane) nhưng tách hai mặt phẳng điều khiển (control plane) cho hai chiến lược:

\begin{itemize}[leftmargin=*]
    \item \textbf{Ingest}: \texttt{Ingestor} gán $T_{event}$, băm $\text{PartitionID}=\text{hash}(host)\bmod 12$, produce vào Kafka topic \texttt{events} (kèm Punctuation Token cho Strict mode).
    \item \textbf{Data Plane (Kafka)}: 12 phân vùng, RF=3, append-only; sắp xếp out-of-order đẩy xuống Worker.
    \item \textbf{Compute (4 Worker)}: ánh xạ 3 phân vùng/node, tách luồng nạp (poll + backpressure) và xử lý (drain theo $T_{event}$, dedupe, gom window, RocksDB).
    \item \textbf{Control Plane (gRPC/HTTP)}: Strict dùng Coordinator HA (hỗ trợ Raft hoặc ZooKeeper Lock); Heuristic dùng Aggregator HA (Active--Standby qua khóa ZooKeeper hoặc File Lock).
    \item \textbf{Tiered Storage}: Tier-1 RocksDB (nóng), Tier-2 Shared Volume (checkpoint 10s), Tier-3 MinIO/S3 (cold \& DR).
\end{itemize}

\section{Thiết kế Tuyến Strict Watermark (0\% Loss)}
Strict Path phục vụ nghiệp vụ kiểm toán/đối soát tài chính, ưu tiên tính đúng đắn tuyệt đối; mốc thời gian toàn cục chỉ tiến khi \emph{mọi} phân vùng đang hoạt động đã vượt qua mốc đó. Theo Özsu \& Valduriez, đây là thiết kế \textbf{PC/EC} (chọn Consistency cả khi có Partition lẫn khi Else), chấp nhận chi phí điều phối cao và độ trễ tăng để bảo đảm tính nhất quán mạnh.

\begin{itemize}[leftmargin=*]
    \item \textbf{Định thời dựa trên Punctuation}: Ingestor chèn token kiểm soát mang mốc cam kết $T_{commit}$. Khi phân vùng rảnh (Idle), Ingestor vẫn gửi \emph{Empty Punctuation Token} mỗi 1000ms để dòng thời gian logic toàn cụm không bị đóng băng --- đây là cơ chế \emph{chính} giữ $W_{global}$ luôn tịnh tiến, giải quyết bài toán \emph{không thể phân biệt phân vùng nhàn rỗi với phân vùng tắc nghẽn} (một vấn đề cốt lõi trong hệ phân tán theo Özsu \& Valduriez).
    \item \textbf{Worker-side Bounded Priority Queue}: Min-Heap RAM tối đa 10.000 bản ghi sắp xếp lại theo $T_{event}$; ép đẩy (\texttt{Forced Pop}) khi hàng đợi đầy hoặc phần tử đỉnh heap bị giữ quá $\texttt{max\_wait\_ms}=1000$ms, chống tràn bộ đệm.
    \item \textbf{Hội tụ toàn cục tối thiểu (Min-of-Min)}:
    {\large
    \begin{equation}
    \boxed{W_{global}(t) = \max\Big(W_{global}(t{-}1),\ \min_{\forall P_k \in \text{Active}} LW_i(P_k)\Big)}
    \end{equation}
    }
    Phân vùng phản hồi chậm (STALE) và nhàn rỗi ngầm (IDLE) \emph{vẫn} tham gia $\min()$ để giữ cam kết 0\% loss. Chỉ phân vùng xác nhận lỗi (FAILED) hoặc được Worker đánh dấu nhàn rỗi tường minh (\texttt{is\_temporary\_idle}) sau \texttt{IDLE\_TIMEOUT\_S}$=30.0$s mới bị loại trừ; khi có sự kiện mới, phân vùng tự động quay lại hàm tính --- đây là cơ chế \textbf{Idleness Bypass} chống Straggler.
    \item \textbf{Chống Split-Brain (Fencing)}: Mọi lệnh điều phối kèm cặp $(term, command\_id)$. Worker từ chối lệnh có $term$ nhỏ hơn nhiệm kỳ lớn nhất đã biết, và bỏ qua lệnh trùng $command\_id$ (idempotent).
    \item \textbf{Tái phân bổ cân bằng tải}: Khi \texttt{HEARTBEAT\_TIMEOUT\_S}$=10$s, Coordinator phát hiện Worker chết, tăng \emph{fencing term}, chia \emph{đều} các phân vùng mồ côi cho các Worker sống (ví dụ P9$\to$node0, P10$\to$node1, P11$\to$node2) --- tránh dồn tải gây sập dây chuyền. Khi $\ge 2$ Worker sập, Cascading Failure Protocol điều phối qua Raft log để tái phân bổ toàn bộ phân vùng mồ côi.
    \item \textbf{Strict Failback State Machine}: Bàn giao ngược 5 bước trên Raft log: \texttt{PAUSE} $\to$\allowbreak \texttt{FLUSH\_ACK} $\to$\allowbreak \texttt{KAFKA\_REASSIGN} $\to$\allowbreak \texttt{SEEK\_RESUME} $\to$\allowbreak \texttt{COMPLETE}, qua chu trình trạng thái \texttt{ASSIGNED}$\to$\allowbreak\texttt{REASSIGNING}$\to$\allowbreak\texttt{PAUSED}$\to$\allowbreak\texttt{ASSIGNED}. Bảo đảm tại một thời điểm duy nhất một Worker được consume/ghi một phân vùng (Exactly-Once output). Output được phát song song vào \texttt{strict\_results} và \texttt{audit\_results} (Critical Audit Sink) để consumer kiểm toán replay độc lập.
    \item \textbf{Exactly-Once input}: Hai vòng lọc khi replay --- Watermark Filter ($T_{event}<W_{global}\Rightarrow$ drop) và bảng băm \texttt{seen-id} (TTL) trong RocksDB; Worker tiếp quản \texttt{seek(offset+1)} từ checkpoint Tier-2.
\end{itemize}

\section{Thiết kế Tuyến Heuristic Watermark (Low Latency)}
Heuristic Path loại bỏ điều phối đồng bộ chặn, trao quyền tự trị cục bộ tối đa cho Worker (Özsu \& Valduriez: \textbf{Site Autonomy} giúp giảm phụ thuộc vào điều phối trung tâm, cải thiện hiệu năng và giảm trễ) và bù dữ liệu muộn bất đồng bộ qua DLQ. Đây là thiết kế \textbf{PA/EL}: khi Partition chọn Availability, khi Else chọn Latency, chấp nhận nhất quán sau cùng (Eventual Consistency) để đạt độ trễ cực thấp.

\begin{itemize}[leftmargin=*]
    \item \textbf{Ước lượng phân vị thích ứng DDSketch}: Đo độ trễ thô $lag=t_{arrival}-T_{event}$ qua Sliding Window DDSketch (60 sub-sketch $\times$ 1s), bảo đảm sai số tương đối $\frac{|\hat q-q|}{q}\le\alpha$ ($\alpha=0.01$), bộ nhớ $\le 100$KB/partition (hard cap \texttt{max\_buckets}=1024, \texttt{MAX\_LAG}=1h). DDSketch có tính \emph{cộng gộp} (mergeable): Aggregator có thể gộp sketch từ nhiều phân vùng để ước lượng phân vị toàn cục với sai số giới hạn.
    \item \textbf{Watermark cục bộ thích ứng (đơn điệu)}:
    {\large
    \begin{equation}
    \boxed{W_{h,i}^{P_k}(t) = \max\Big(W_{h,i}^{P_k}(t{-}1),\ T_{event\_max} - L_{eff}(t)\Big),\quad L_{eff}=\text{DDSketch.quantile}(p)}
    \end{equation}
    }
    Phân vị tự nâng từ $p_{normal}=0.99$ lên $p_{safe}=0.999$ khi phát hiện burst lag $>2.0\times$ trung bình, mở rộng biên an toàn (Loss Budget Guarantee).
    \item \textbf{Cold Start \& Negative Lag}: 3 pha khởi động (buffer-only $\to$ Conservative Prior $L_{max}=60$s $\to$ Normal khi $\ge 10$s \emph{và} $\ge 1000$ mẫu); lọc lag âm do lệch đồng hồ (Clock Skew), hiệu chuẩn cộng bù $|\text{median\_neg\_lag}|$ khi tỷ lệ âm $\ge 1\%$.
    \item \textbf{Aggregator HA (Active--Standby)}: Sử dụng khóa loại trừ trên ZooKeeper tại \texttt{/csdlpt/aggregator-lock} hoặc File Lock dự phòng (\texttt{FileLockLeader} dùng \texttt{fcntl}/\texttt{msvcrt}). Active ghi file nhịp tim mỗi 1s. Nếu nhịp tim mất hoặc quá hạn ($>1.5$s), Standby tự động chiếm khóa, khôi phục trạng thái từ RocksDB/JSON và lên Active ($RTO\le 2$s), tránh SPOF.
    \item \textbf{Replay-Mode Sketch Snapshot \& Rollback}: Hàng đợi 6 snapshot DDSketch trong 60s gần nhất. Khi phát hiện trễ tăng $10\times$ (dấu hiệu replay), Worker rollback về snapshot an toàn và đóng băng bộ ước lượng, tránh ``nhiễm bẩn'' thống kê. Kết hợp Replay Sub-Checkpointing để không phải replay lại từ đầu khi xử lý backlog lớn.
    \item \textbf{Đường ống sửa lỗi DLQ Correction}: Bản ghi $T_{event}<W_{global\_h}$ định tuyến vào \texttt{late\_logs\_dlq} (không drop ngầm, hạch toán \texttt{per\_window\_loss}). DLQ Consumer nạp lại trạng thái cửa sổ lịch sử, tính Delta và phát \texttt{CorrectionMessage} xuống downstream (Incremental / Replace / Append-with-Versioning) $\Rightarrow$ Eventual Completeness 100\%. Mô hình này hiện thực hóa \textbf{Eventual Consistency} theo Özsu \& Valduriez: chấp nhận nhất quán tức thời yếu hơn để giảm trễ, bù đắp bằng sửa lỗi bất đồng bộ.
\end{itemize}

\section{Quản lý Trạng thái \& Vòng đời Dữ liệu}
\textbf{Máy trạng thái sở hữu phân vùng (Partition Ownership FSM) --- dùng chung cho cả hai tuyến:} Mỗi phân vùng trải qua vòng đời xác định, được duy trì trên Raft log (Strict) hoặc Aggregator state (Heuristic):

\begin{center}
\small
\setlength{\tabcolsep}{3pt}
\begin{tabular}{@{}>{\raggedright\arraybackslash}p{2.0cm}>{\raggedright\arraybackslash}p{2.7cm}>{\raggedright\arraybackslash}p{5.2cm}>{\raggedright\arraybackslash}p{5.2cm}@{}}
\toprule
\textbf{Trạng thái} & \textbf{Trigger} & \textbf{Hành vi} & \textbf{Điểm khác biệt Strict vs Heuristic} \\
\midrule
\texttt{ASSIGNED} & Khởi tạo / Failback hoàn tất & Worker consume+ghi bình thường; heartbeat báo cáo $LW_i$(P$_k$) mỗi 1s. & Strict: Coordinator giám sát heartbeat, tính $W_{global}$. Heuristic: Worker tự chốt cửa sổ, gửi $W_{h}$ lên Aggregator mỗi 500ms. \\
\midrule
\texttt{REASSIGNING} & Heartbeat timeout (10s) / Worker sập & Coordinator/Aggregator tăng fencing term, chọn Worker mới, gửi lệnh PAUSE tới Worker cũ. & Strict: quyết định tái phân bổ ghi vào Raft log (quorum acks). Heuristic: Aggregator cập nhật bảng ánh xạ, báo qua ZooKeeper. \\
\midrule
\texttt{PAUSED} & Nhận lệnh PAUSE từ Coordinator & Worker dừng consume, flush RocksDB + Kafka offset lên Tier-2, gửi FLUSH\_ACK. & Strict: bắt buộc qua 5 bước failback trên Raft log. Heuristic: snapshot DDSketch hiện tại, đóng băng ước lượng. \\
\midrule
\texttt{ORPHANED} & $\ge 2$ Worker sập đồng thời & Cascading Failure Protocol: Coordinator phân phối đều phân vùng mồ côi cho Worker sống, Worker mới tải checkpoint Tier-2, \texttt{seek(offset+1)}. & Strict: tái phân bổ ghi vào Raft log, fencing term mới. Heuristic: Aggregator tự phân bổ lại, Worker vào Cold Start (xây lại DDSketch từ đầu). \\
\bottomrule
\end{tabular}
\end{center}

\vspace{4pt}
\noindent\textbf{Luồng dữ liệu end-to-end --- truy vết một sự kiện log:}
\begin{enumerate}[leftmargin=*,nosep,topsep=2pt]
    \item \textbf{Ingest}: \texttt{Ingestor} đọc CSV, gán $T_{event}=\text{parse\_ts}(row)$, tính $\text{PartitionID}=\text{hash}(host)\bmod 12$, produce vào Kafka topic \texttt{events}. Nếu là Strict mode: chèn thêm Punctuation Token mang $T_{commit}$ (hoặc Empty Token mỗi 1000ms nếu phân vùng rảnh).
    \item \textbf{Kafka Buffering}: Bản ghi nằm trong partition log với offset đơn điệu tăng. RF=3 bảo đảm không mất dữ liệu nếu 1 broker chết.
    \item \textbf{Worker Poll \& Backpressure}: Worker poll batch từ Kafka. Nếu hàng đợi nội bộ $\ge 500$: \texttt{consumer.pause()}; vơi còn $\le 100$: \texttt{resume()}.
    \item \textbf{Reorder (Bounded PQ)}: Đẩy vào Min-Heap (max 10.000 phần tử) theo $T_{event}$. Drain đỉnh heap khi $T_{event} \le W_{global}$ (Strict) hoặc $T_{event} - \text{DDSketch.quantile}(p) \le T_{event\_max}$ (Heuristic). Forced Pop nếu heap đầy hoặc phần tử đỉnh bị giữ quá 1000ms.
    \item \textbf{Deduplication}: Tra bảng băm \texttt{seen-id} (TTL) trong RocksDB. Nếu đã tồn tại $\Rightarrow$ drop (idempotent replay).
    \item \textbf{Window Aggregation}: Gom bản ghi vào Tumbling Window 5s theo $T_{event}$. Khi watermark vượt biên phải cửa sổ $\Rightarrow$ chốt, tính toán aggregate, gán \texttt{window\_id} duy nhất.
    \item \textbf{Phân kỳ tại điểm chốt cửa sổ}:
    \begin{itemize}[leftmargin=*,nosep]
        \item \textbf{Strict}: Phát kết quả vào \texttt{strict\_results} và \texttt{audit\_results} (dual sink). Ghi checkpoint RocksDB + offset lên Tier-2 mỗi 10s. \emph{Điểm chốt là bất biến --- không bao giờ sửa lại.}
        \item \textbf{Heuristic}: Phát kết quả vào output topic. Bản ghi $T_{event} < W_{global\_h}$ (đến muộn) $\Rightarrow$ định tuyến vào \texttt{late\_logs\_dlq}. DLQ Consumer nạp lại trạng thái cửa sổ lịch sử, tính Delta, phát \texttt{CorrectionMessage}. \emph{Điểm chốt có thể được sửa đổi bất đồng bộ đến khi đạt Eventual Completeness 100\%.}
    \end{itemize}
    \item \textbf{Tiered Eviction}: Cửa sổ đã chốt $\to$ \texttt{CLOSED}. Định kỳ: \texttt{UPLOADING}\allowbreak\ (đẩy lên MinIO, key \texttt{\{partition\}/\{window\_id\}}), \texttt{UPLOADED}\allowbreak\ (xác nhận tồn tại trên S3), \texttt{PURGED}\allowbreak\ (xóa RocksDB cục bộ). Key deterministic chống trùng lặp khi re-upload sau failover.
\end{enumerate}

\vspace{2pt}
\noindent\textbf{Tóm lược phân kỳ PACELC:} Strict chọn \textbf{PC/EC} --- khi Partition dừng watermark (Consistency), khi Else chờ đủ Punctuation (Consistency); Heuristic chọn \textbf{PA/EL} --- khi Partition Worker vẫn chốt cửa sổ (Availability), khi Else chốt sớm qua DDSketch (Latency), bù đắp bằng DLQ Correction sau.

\end{document}
